% Introduction and Related Work for PRIZM (zero-shot skeleton action recognition).
% Use: \input{intro_related.tex} inside main.tex, after \maketitle.
% Needs a bibliography; \cite keys used below:
%   Skeleton ZSL : reVISE17 jpose19 cadavae19 synse21 smie23 gzssar23 purls24 savae24 star24
%                  starpp26 dvta25 infocpl25 scople25 neuron25 fsvae25 tdsm25 flora26
%                  skelcache dynapurls skeletoncontext amsf25 pgfa25
%   Skeleton data: ntu60 ntu120 pku shiftgcn
%   Generic ZSL  : lampert14 xian18 wang18kg koh20cbm clip21
%   NeSy         : logiczsl25 andolfi25 pcbear
%
% Standalone: compile this file directly (Overleaf: Menu -> Main document = intro_related.tex),
% with refs.bib uploaded in the same folder. Compiler: pdfLaTeX.

\documentclass[11pt]{article}
\usepackage[margin=1in]{geometry}
\usepackage[T1]{fontenc}
\usepackage[utf8]{inputenc}
\usepackage{amsmath,amssymb}
\usepackage{booktabs}
\usepackage{cite}
\usepackage{graphicx}
\usepackage{multirow}
\usepackage[table]{xcolor}
\usepackage[hidelinks]{hyperref}

\title{PRIZM: Neuro-Symbolic Primitive-Rule Priors for Zero-Shot Skeleton Action Recognition}
\author{}
\date{}

\begin{document}
\maketitle

\section{Introduction}

Human action recognition is a core problem in video understanding, with applications in human-computer interaction, healthcare monitoring, sports analysis, and surveillance. Among the available modalities, 3D skeleton sequences have received growing attention. They can be obtained reliably from depth sensors or pose estimators, they are compact and privacy-preserving, and they are robust to changes in appearance, illumination, viewpoint, and background~\cite{ntu60,ntu120}. Driven by graph convolutional networks~\cite{stgcn18,agcn19,shiftgcn,ctrgcn21}, supervised skeleton-based models have reached high accuracy on standard benchmarks. These models, however, can only recognize the categories they were trained on. Human actions form an open-ended set, new categories emerge with every new application, and collecting and annotating skeleton data for each of them is expensive. This has motivated \emph{zero-shot} skeleton action recognition (ZSAR), in which a model is trained on a set of seen classes and must recognize unseen classes using only their semantic descriptions~\cite{jpose19,synse21}. In the conventional ZSL setting the test samples come only from unseen classes, while in the generalized zero-shot setting, referred to as GZSL, they come from both seen and unseen classes.

The central difficulty of ZSAR is that the model must transfer knowledge from seen to unseen classes through the semantic space, without ever observing a skeleton of an unseen class. Most existing methods follow an \emph{align-then-classify} paradigm. A skeleton encoder and a language encoder, typically a pretrained vision-language model such as CLIP~\cite{clip21}, are aligned on the seen classes, and an unseen sample is assigned to the class whose semantic embedding is closest to its skeleton feature~\cite{synse21,smie23,star24,neuron25}. Research along this line has developed in two directions. The first improves the alignment and the classifier, through generative and variational models~\cite{cadavae19,savae24,fsvae25}, diffusion~\cite{tdsm25}, prompt learning~\cite{star24,scople25}, prototype modeling~\cite{neuron25}, and distribution-aware classifiers~\cite{flora26}. The second enriches the class semantics themselves, replacing bare class names with descriptions generated by large language models (LLMs), with part-level and phase-level descriptions, or with contextual cues about the objects and scenes involved~\cite{purls24,infocpl25,dynapurls,skeletoncontext}. Both directions have led to steady progress on the NTU RGB+D benchmarks~\cite{ntu60,ntu120}.

Despite this progress, the semantic side of ZSAR still represents each class as free-form text that is mapped to a vector by a language encoder. We argue that this representation has three limitations for skeleton data. First, the geometry of the semantic space is inherited from language rather than from motion. Two actions may differ in a detail that matters for the body but hardly changes the sentence. A typical case is a pair of actions that are temporal reversals of each other, such as \emph{wearing} and \emph{taking off} a jacket: both involve the same body parts and the same objects, and differ mainly in the order of their sub-motions. Their text embeddings are therefore nearly identical, and the classifier receives little semantic signal to separate them. Second, skeletons do not contain the appearance cues, such as objects and scenes, that make many actions easy to distinguish in RGB videos. Descriptions that refer to objects, such as a pen or a keyboard, are thus only weakly grounded in the skeleton, and what remains observable is body motion and its temporal organization. Third, free-form descriptions do not state explicitly \emph{what} a seen and an unseen class have in common, for instance that both hands are active, that one hand moves toward the head, or that the body bends forward. Knowledge transfer therefore depends on how the language model happens to place the two descriptions, and it can be neither inspected nor corrected when it fails.

A classical way to make shared structure explicit is to describe classes by attributes or primitives that are reused across categories~\cite{lampert14,xian18}. Human actions are well suited to this view. A limited vocabulary of body-part motions, spatial relations between body parts, interactions between two persons, and temporal orderings is enough to describe a large number of actions, and most of these primitives can be observed directly in a skeleton sequence. Such a vocabulary can be organized as a knowledge graph (KG) that links each action to the primitives it involves. The same primitives appear in seen and unseen classes, so the KG states explicitly the shared structure on which transfer should rely. Moreover, the KG is written in terms of motion and can therefore describe differences that the text does not capture.

Combining a KG with a modern ZSAR model, however, raises two questions. The first is how to represent constraints that free-form text handles poorly, in particular temporal order and negation, that is, which sub-motion comes first and which motion must \emph{not} occur. The second is how to inject the KG without introducing components that overfit the seen classes. Learned modules over symbolic structures are known to exploit \emph{reasoning shortcuts}, in which the model satisfies the training labels with concepts that do not have their intended meaning~\cite{marconato26,andolfi25}. In zero-shot learning this risk is amplified: the semantic module is trained only on the seen classes, and the unseen classes provide no labels against which the learned concepts could be checked. A KG module that is flexible enough to fit the seen classes may therefore distort exactly the knowledge that should transfer to the unseen ones.

In this paper we propose \textbf{PRIZM} (PRImitive-rule priors for Zero-shot skeleton Motion recognition), a neuro-symbolic framework that augments the class semantics of a neural ZSAR model with priors derived from a symbolic primitive-rule KG. In the KG, each action class is linked to the motion primitives it \emph{requires}, and actions that are temporal reversals of each other are additionally linked to ordered primitives that they \emph{forbid}. PRIZM turns these rules into class embeddings and fuses them with the text semantics of the underlying model. Following the observation above, all KG-related operations are fixed and contain no learnable parameters. The KG thus acts as a prior on the semantic space, while all learning remains in the skeleton-text model. We instantiate PRIZM on Flora~\cite{flora26}, a recent model that combines neighbor-aware semantic attunement with a flow-matching classifier, and evaluate it on the NTU RGB+D 60 and 120 benchmarks. In addition to the standard comparisons, we design controlled experiments to verify that the improvements come from the content of the KG rather than from the additional semantic term itself.

The main contributions of this paper are summarized as follows:
\begin{itemize}
  \item We propose PRIZM, a framework that injects a primitive-rule knowledge graph into the class semantics of zero-shot skeleton action recognition through fixed, parameter-free operations, which makes the transferred knowledge explicit and inspectable while avoiding overfitting to the seen classes.
  \item We introduce Order-aware Rule Expansion (ORE) and Polarity-aware Primitive Composition (PPC), which express temporal order and negation in the class semantics, so that actions that differ only in the order of their sub-motions obtain distinct representations.
  \item We design COA-Fuse, a parameter-free fusion of KG semantics and text semantics that keeps the complementary part of the KG and adapts its contribution to each class.
  \item We conduct extensive experiments on the NTU RGB+D 60 and NTU RGB+D 120 benchmarks, together with controlled analyses, including a shuffled-KG control and comparisons with learned KG modules, which verify that the improvements come from the content of the knowledge graph.
\end{itemize}

The remainder of this paper is organized as follows. Section~\ref{sec:related} reviews related work. Section~\ref{sec:method} describes the proposed method. Section~4 presents the experimental results and analyses, and Section~5 concludes the paper.


\section{Related Work}
\label{sec:related}

\subsection{Zero-Shot Skeleton-Based Action Recognition}

ZSAR methods typically use a skeleton backbone pretrained on the seen classes, most commonly Shift-GCN~\cite{shiftgcn}, and focus on relating its features to the semantics of unseen classes. The task was established by SynSE~\cite{synse21}, which adapted ReViSE~\cite{reVISE17} and JPoSE~\cite{jpose19} as baselines and learned generative embeddings guided by the parts of speech of the class name. Early methods align the two modalities in a shared space with variational autoencoders~\cite{cadavae19,gzssar23} or by maximizing their mutual information~\cite{smie23}. Later methods align them at a finer granularity: PURLS~\cite{purls24} aligns body parts with part-level descriptions, STAR~\cite{star24,starpp26} learns prompts guided by body-part side information, SA-DVAE~\cite{savae24} aligns only the semantic-related part of the skeleton features, and SCoPLe~\cite{scople25} learns cross-modal prompts for both encoders. Recent works have also revisited the classifier. Neuron~\cite{neuron25} builds context-aware micro-prototypes that evolve over space and time, FS-VAE~\cite{fsvae25} enhances the low-frequency components of joint trajectories, and TDSM~\cite{tdsm25} aligns skeleton and text with a diffusion model. Flora~\cite{flora26}, on which we build, attunes each class semantic with its neighboring classes and replaces the nearest-neighbor decision with a flow-matching classifier. For the generalized setting, the bias toward seen classes is reduced by calibrated stacking~\cite{chao16} or gating~\cite{gzssar23,scople25}, and test-time adaptation methods~\cite{skelcache,dynapurls} further exploit the unlabeled test stream in a transductive setting. These methods mainly improve the alignment, the classifier, or the adaptation procedure, while the class semantics remain text embeddings. Our work is orthogonal: it changes the class semantics by adding a structured prior, and it could in principle be combined with any of these models.

\subsection{Semantic Enrichment}

Class names provide limited and often ambiguous information, and several works therefore enrich them with additional text. LLM-generated descriptions of how an action is performed are now widely used~\cite{gzssar23,purls24,star24}. InfoCPL~\cite{infocpl25} generates many descriptions per class and selects the most relevant ones for each sample, AMSF~\cite{amsf25} fuses several types of expert-written and LLM-generated descriptions with attention, and part-level and phase-level descriptions describe individual body parts and temporal phases separately~\cite{purls24,dynapurls}. SkeletonContext~\cite{skeletoncontext} adds context about the environment and the objects involved, and GenPrior~\cite{genprior26} derives kinematic prototypes from a text-to-motion model. These methods improve the descriptive quality of the semantics, but each description remains a free-form sentence encoded by a language model, so the relations between classes are determined implicitly by the language embedding space. In contrast, our KG states explicitly which primitives a class shares with other classes, and represents temporal order and negation through ordered primitives and forbid rules.

\subsection{Knowledge Graphs and Neuro-Symbolic Learning}

Attribute-based zero-shot learning recognizes unseen classes from a shared vocabulary of human-defined attributes~\cite{lampert14}, and knowledge graphs have been used to propagate classifier weights between related categories with graph convolutional networks~\cite{wang18kg,dgp19}. Concept bottleneck models predict interpretable concepts before the class~\cite{koh20cbm}, including pose-based concepts for action recognition~\cite{pcbear}, but they are designed for the supervised setting. Neuro-symbolic methods integrate symbolic knowledge into neural models through differentiable losses or logic programs~\cite{xu18semloss,deepproblog18,ltn22}. LogicZSL~\cite{logiczsl25} converts LLM-generated rules into fuzzy losses for compositional zero-shot learning, but its rules express only implications, without negation or temporal order. When concepts are not supervised directly, such models are prone to reasoning shortcuts~\cite{marconato26}, which prototype-based grounding has been proposed to mitigate~\cite{andolfi25}. PRIZM avoids this problem by design: the KG is used only through fixed composition and fusion operations, so that the symbolic knowledge shapes the semantic space without being reinterpreted to fit the seen classes.

\section{Method}
\label{sec:method}

This section first formulates the task and summarizes the Flora baseline. It then describes the primitive-rule knowledge graph, its order-aware extension ORE, the PPC module that turns rules into class embeddings, and the COA-Fuse procedure that combines these embeddings with the text semantics. An overview of PRIZM is shown in Fig.~\ref{fig:overview}.

\begin{figure}[t]
  \centering
  \includegraphics[width=\linewidth]{prizm.jpg}
  \caption{Overview of PRIZM. The primitive-rule KG is extended with ordered primitives and reverse-order constraints by ORE. PPC turns the rules of each class into a KG embedding, and COA-Fuse adds this embedding to the attuned text semantics of Flora. All KG operations are fixed and have no learnable parameters.}
  \label{fig:overview}
\end{figure}

\subsection{Problem Formulation}

Let $\mathcal{Y} = \mathcal{Y}^s \cup \mathcal{Y}^u$ be the set of action classes, where the seen classes $\mathcal{Y}^s$ and the unseen classes $\mathcal{Y}^u$ are disjoint, and let $n = |\mathcal{Y}|$. The training set $\mathcal{D}^{tr} = \{(\mathbf{x}_i, y_i)\}$ contains skeleton sequences $\mathbf{x}_i$ with labels $y_i \in \mathcal{Y}^s$ only. Every class $c \in \mathcal{Y}$, seen or unseen, is described by side information, which in our case consists of the class name, an LLM-generated description, and a set of rules in a knowledge graph. In the ZSL setting, a test sample must be assigned to a class in $\mathcal{Y}^u$; in the GZSL setting, the label space is $\mathcal{Y}$. As in prior work~\cite{synse21,star24,flora26}, the skeleton encoder is a Shift-GCN~\cite{shiftgcn} pretrained on the seen classes and kept frozen, and each sequence is represented by its pooled feature $\mathbf{f} \in \mathbb{R}^{256}$.

\subsection{Preliminaries: Flora}
\label{sec:flora}

\paragraph{Text semantics.}
Flora~\cite{flora26} encodes each class with the frozen text encoder of CLIP~\cite{clip21}. For a sentence, the encoder outputs $L$ token embeddings and one pooled embedding, which are stacked into a matrix in $\mathbb{R}^{(L+1) \times 768}$. Each class is described by two sentences, the class name and an LLM-generated description, and the two matrices are concatenated along the feature dimension, giving the class semantics $\mathbf{T}_c \in \mathbb{R}^{(L+1) \times 1536}$.

\paragraph{Neighbor-aware attunement.}
Let $s_{cj}$ be the cosine similarity between the token-averaged semantics of classes $c$ and $j$, and let $\mathcal{N}_k(c)$ be the $k$ classes most similar to $c$. Flora attunes each class with its neighbors,
\begin{equation}
  \mathbf{A}_c = \mathbf{T}_c + \frac{\omega}{k} \sum_{j \in \mathcal{N}_k(c)} s_{cj}\, \mathbf{T}_j ,
  \label{eq:attune}
\end{equation}
where $\omega$ is a mixing weight. The attuned semantics $\mathbf{A}_c$ are computed once for all classes and are used both for training on the seen classes and for recognition of the unseen ones.

\paragraph{Alignment and flow-matching classifier.}
Training has two phases. In the first phase, a semantic encoder and a motion encoder map $\mathbf{A}_{y}$ and $\mathbf{f}$ to Gaussian latents $\mathbf{z}_0$ and $\mathbf{z}_1$, and are trained to match the means and variances of the two latents together with within-modality and cross-modality reconstruction losses. In the second phase, the encoders are frozen and a DiT velocity predictor $v_\theta$ is trained with flow matching to transport the semantic latent to the skeleton latent along the path $\mathbf{z}_t = (1 - (1-\sigma)t)\,\mathbf{z}_0 + t\,\mathbf{z}_1$, whose target velocity is $\mathbf{u} = \mathbf{z}_1 - (1-\sigma)\mathbf{z}_0$. A contrastive term pushes the predicted velocity away from the target velocity of another training pair. At test time, for every candidate class $c$ the latent $\mathbf{z}_0^{c}$ of $\mathbf{A}_c$ is paired with the latent $\mathbf{z}_1$ of the test sample, and the sample is assigned to the class with the lowest velocity error at a fixed time $t^\ast$,
\begin{equation}
  \hat{y} = \arg\min_{c} \big\| v_\theta(\mathbf{z}^{c}_{t^\ast}, t^\ast) - \mathbf{u}^{c} \big\|_2^2 .
  \label{eq:flora-pred}
\end{equation}
PRIZM leaves the encoders, the velocity predictor, the losses, and the decision rule of Flora unchanged. It only replaces the class semantics $\mathbf{A}_c$ in Eq.~\eqref{eq:flora-pred} and in training by the KG-augmented semantics $\mathbf{S}_c$ defined in Section~\ref{sec:coa}.

\subsection{Primitive-Rule Knowledge Graph}
\label{sec:kg}

\paragraph{Primitive vocabulary.}
The KG is built on a vocabulary $\mathcal{P}$ of 82 motion primitives that describe observable properties of a skeleton sequence. The primitives cover seven body groups, namely hand, two-person interaction, whole body, arm, torso, leg, and head, and three types: \emph{states}, such as a hand being at the face; \emph{motions}, such as a fast forward thrust of the hand; and \emph{relations}, such as the hands of two persons touching. Each primitive $p$ has a short natural-language description $d_p$, for example ``A hand is at the front of the head.''

\paragraph{Class rules.}
Each class $c$ is linked to the primitives through three relations. The set $\mathcal{R}_c \subseteq \mathcal{P}$ contains the primitives that must occur at some time during the action, the set $\mathcal{F}_c \subseteq \mathcal{P}$ contains the primitives that must not occur, and the sequence $\mathbf{q}_c = (q_1, \dots, q_m)$ lists required primitives in the order in which they occur. Primitives not mentioned in the rules of a class are irrelevant to it. For example, \emph{drink water} requires a hand at the face, a hand at head level, and a flexed arm, and forbids, among others, a second person and a hand above the head. On the 120 classes of NTU RGB+D 120, a class requires 3.1 primitives and forbids 2.6 primitives on average, and 27 classes have a non-empty sequence.

\paragraph{Construction.}
The rules were generated by an LLM from the class names and the primitive vocabulary only. No skeleton data and no experimental result were used, so the KG contains the same kind of class-level knowledge as the class names and descriptions used by Flora, and is available for unseen classes without violating the zero-shot protocol. The rules were then revised once by a rule-level analysis that detects classes with identical rule sets and class pairs separated by at most two primitives.

\subsection{Order-aware Rule Expansion}
\label{sec:seqkg}

Several pairs of NTU classes are temporal reversals of each other, such as \emph{wear jacket} and \emph{take off jacket}, \emph{put on a shoe} and \emph{take off a shoe}, or \emph{put something into a bag} and \emph{take something out of a bag}. The two classes of such a pair involve the same primitives in reverse order. Their sets $\mathcal{R}_c$ and $\mathcal{F}_c$ are therefore identical, and any representation built from these sets alone cannot separate them. In the original KG, seven class pairs are indistinguishable for this reason.

Order-aware Rule Expansion, abbreviated as ORE, makes the order explicit without adding new knowledge, using only the sequences $\mathbf{q}_c$ already present in the KG. For every pair of consecutive primitives $(a, b)$ in the sequence of some class, an \emph{ordered primitive} $a \!\succ\! b$ is added to the vocabulary, with the description
\begin{equation}
  d_{a \succ b} = \text{``First, } d_a \text{ Then, } d_b \text{''}.
  \label{eq:seqdesc}
\end{equation}
For each class $c$, the ordered primitives of its own sequence are added to $\mathcal{R}_c$. If the reverse primitive $b \!\succ\! a$ also exists, because another class performs the same sub-motions in the opposite order, it is added to $\mathcal{F}_c$. On NTU RGB+D 120, this procedure creates 29 ordered primitives and 14 reverse-order constraints, and the number of class pairs with identical rules drops from seven to zero. The extended vocabulary $\mathcal{P}'$ has 111 primitives.

\subsection{Polarity-aware Primitive Composition}
\label{sec:compose}

Each primitive description is encoded with the same frozen CLIP text encoder as the class texts, giving $\mathbf{E}_p \in \mathbb{R}^{(L+1) \times 768}$. A class embedding is then obtained as a weighted combination of the primitive embeddings,
\begin{equation}
  \mathbf{K}_c = \sum_{p \in \mathcal{P}'} w_{cp}\, \mathbf{E}_p .
  \label{eq:compose}
\end{equation}
A plain composition uses only the required primitives, with uniform weights $w_{cp} = r_{cp} = \mathbb{1}[p \in \mathcal{R}_c] / |\mathcal{R}_c|$. This ignores $\mathcal{F}_c$, so the reverse-order constraints of ORE would have no effect, and the two classes of a reversed pair would still differ only by their own ordered primitive. Polarity-aware Primitive Composition, abbreviated as PPC, additionally subtracts the forbidden ordered primitives. Let $\mathcal{F}^{o}_c \subseteq \mathcal{F}_c$ be the forbidden ordered primitives of class $c$, and let $f_{cp} = \mathbb{1}[p \in \mathcal{F}^{o}_c] / |\mathcal{F}^{o}_c|$. The weights are
\begin{equation}
  w_{cp} =
  \begin{cases}
    (1+\lambda)\, r_{cp} - \lambda\, f_{cp}, & \text{if } \mathcal{F}^{o}_c \neq \emptyset, \\
    r_{cp}, & \text{otherwise},
  \end{cases}
  \label{eq:signed}
\end{equation}
where $\lambda \ge 0$ controls the strength of the constraint. The weights of each class sum to one, so the component shared by all CLIP embeddings and the overall scale of $\mathbf{K}_c$ are preserved, and the weight of the KG term does not need to be retuned. Eq.~\eqref{eq:signed} can be written as $\mathbf{r}_c + \lambda(\mathbf{r}_c - \mathbf{f}_c)$: the composition moves away from the reverse order along the direction that separates the two orders.

Only forbidden \emph{ordered} primitives are used in Eq.~\eqref{eq:signed}. Ordinary forbid rules, such as the absence of a second person, apply to most classes and describe what the action is not, rather than what distinguishes it from a similar class. Including them would change the embedding of almost every class. The ordered constraints, in contrast, act only on the 14 classes that belong to a reversed pair.

The composed embedding has 768 dimensions per token, whereas the text semantics of Flora have 1536. We therefore repeat $\mathbf{K}_c$ along the feature dimension, so that it is added to both the class-name half and the description half of $\mathbf{A}_c$.

\subsection{COA-Fuse}
\label{sec:coa}

The simplest way to use $\mathbf{K}_c$ is to add it to the attuned semantics, $\mathbf{S}_c = \mathbf{A}_c + \beta \mathbf{K}_c$, with a global weight $\beta$. This raw addition has three drawbacks. First, CLIP text embeddings are strongly anisotropic, so a large part of $\mathbf{K}_c$ is a component common to all classes, which shifts every class in the same way and does not help to discriminate between them. Second, part of $\mathbf{K}_c$ repeats information already present in the class text. Third, the KG is equally useful for every class, although it is most informative where the text is least reliable. COA-Fuse addresses these three points with three parameter-free steps.

\paragraph{Center.}
The class mean of the KG term is removed,
\begin{equation}
  \tilde{\mathbf{K}}_c = \mathbf{K}_c - \frac{1}{n} \sum_{j \in \mathcal{Y}} \mathbf{K}_j ,
  \label{eq:center}
\end{equation}
so that only the part that differs between classes is added to the semantics. The mean is taken over all seen and unseen classes, using only their rules.

\paragraph{Orthogonalize.}
Let $\tilde{\mathbf{A}}_c = \mathbf{A}_c - \frac{1}{n}\sum_{j} \mathbf{A}_j$ be the centered text semantics. For each token $\ell$, the component of the KG term along the text of the same class is removed,
\begin{equation}
  \mathbf{K}^{\perp}_{c,\ell} = \tilde{\mathbf{K}}_{c,\ell} - \frac{\langle \tilde{\mathbf{K}}_{c,\ell}, \tilde{\mathbf{A}}_{c,\ell} \rangle}{\|\tilde{\mathbf{A}}_{c,\ell}\|_2^2}\, \tilde{\mathbf{A}}_{c,\ell} ,
  \label{eq:orth}
\end{equation}
and the result is rescaled to the original magnitude, $\hat{\mathbf{K}}_c = \mathbf{K}^{\perp}_c \cdot \|\tilde{\mathbf{K}}_c\|_F / \|\mathbf{K}^{\perp}_c\|_F$. This step changes only the direction of the KG term, so that the KG contributes information that the text of the class does not already contain.

\paragraph{Adaptive weight.}
The KG is most useful for classes that the text considers similar to another class while the rules consider them different. Let $\mathbf{S}^{txt}$ be the text similarity matrix used by Flora, and let $\mathbf{S}^{kg}$ be the cosine similarity between the rule vectors of the classes, in which a required primitive has weight $1$ and a forbidden one weight $-0.5$. Both matrices are standardized by the mean and standard deviation of their off-diagonal entries, denoted $z(\cdot)$, since they have different scales. The disagreement of class $c$ is
\begin{equation}
  \delta_c = \max_{j \neq c} \big[ z(S^{txt}_{cj}) - z(S^{kg}_{cj}) \big]_+ ,
  \label{eq:disagree}
\end{equation}
where $[\cdot]_+$ denotes the positive part. The class-specific weight is
\begin{equation}
  \beta_c = \beta \cdot \min\!\Big( \frac{\delta_c}{\bar{\delta}},\, \gamma \Big),
  \label{eq:adaptive}
\end{equation}
where $\bar{\delta}$ is the mean disagreement over all classes and $\gamma$ caps the weight. A class whose text and rules agree receives a small weight, and a class whose text is confused with another class that the rules separate receives a large one.

\paragraph{Final semantics.}
The KG-augmented semantics are
\begin{equation}
  \mathbf{S}_c = \mathbf{A}_c + \beta_c\, \hat{\mathbf{K}}_c ,
  \label{eq:final}
\end{equation}
and replace $\mathbf{A}_c$ everywhere in Flora. The neighbor selection in Eq.~\eqref{eq:attune} is computed from the text similarity alone and is not affected by the KG.

\subsection{Training and Inference}
\label{sec:train}

PRIZM computes $\mathbf{S}_c$ once, before training, for all seen and unseen classes. All steps in Sections~\ref{sec:seqkg}--\ref{sec:coa} are closed-form operations on frozen CLIP embeddings and on the rules. PRIZM therefore adds no learnable parameters to Flora and does not change its losses or its optimization. During training, the alignment and flow-matching objectives of Flora are applied with $\mathbf{S}_y$ in place of $\mathbf{A}_y$ for the seen classes. At test time, the prediction is given by Eq.~\eqref{eq:flora-pred} with $\mathbf{S}_c$ in place of $\mathbf{A}_c$ for the candidate classes. The additional cost consists of encoding the 111 primitive descriptions once and of a few matrix operations on $n$ class embeddings, which is negligible compared with training Flora.

Keeping the KG fixed is a deliberate design choice. A learned module over the KG, such as a graph network or a gate, is trained only on the seen classes and can fit their labels in ways that do not reflect the meaning of the primitives~\cite{marconato26,andolfi25}, and nothing in the training signal checks whether the learned KG representation is still correct for the unseen classes. A fixed composition, in contrast, encodes the same rules in the same way for seen and unseen classes, so that the structure stated in the KG is transferred unchanged. We compare the two options empirically in Section~4.

\section{Experiments}
\label{sec:exp}

\subsection{Experimental Settings}

\paragraph{Datasets and metrics.}
Following Flora~\cite{flora26}, we evaluate PRIZM on two 3D skeleton action recognition benchmarks under the Cross-Subject (Xsub) protocol: \textbf{NTU RGB+D 60}~\cite{ntu60} (56,880 sequences, 60 classes) with the 55/5 and 48/12 seen/unseen splits, and \textbf{NTU RGB+D 120}~\cite{ntu120} (114,480 sequences, 120 classes) with the 110/10 and 96/24 splits, all taken from~\cite{synse21}. In the ZSL setting, we report the top-1 accuracy on the unseen test samples. In the GZSL setting, the label space contains all classes, and we report the accuracy on the seen test samples ($\mathcal{S}$), the accuracy on the unseen test samples ($\mathcal{U}$), and their harmonic mean $\mathcal{H} = 2\mathcal{S}\mathcal{U}/(\mathcal{S} + \mathcal{U})$.

\paragraph{Implementation details.}
We use the released configurations of Flora and keep its architecture, losses, and training schedule unchanged. Each skeleton sequence is represented by a frozen 256-D pooled feature of a Shift-GCN~\cite{shiftgcn} pretrained on the seen classes of the split. As in Flora, we report results with the 1s-Shift-GCN features released by STAR~\cite{star24} and with the 4s-Shift-GCN features released by SynSE~\cite{synse21}, the latter marked with $\dagger$. Class names, LLM descriptions, and primitive descriptions are encoded with the frozen text encoder of CLIP ViT-L/14@336px~\cite{clip21}. The semantic and skeleton VAEs use two-layer MLP encoders and decoders, and the velocity predictor is a single-layer DiT trained with logit-normal timestep sampling; at test time the velocity error is evaluated at $t^\ast = 0.5$. Training consists of 1,000 iterations of alignment and 200 iterations of flow matching, with learning rate $10^{-4}$, weight decay $0.01$, batch size $64$, and $\lambda_{\text{Align}} = \lambda_{\text{Flow}} = 0.1$. In GZSL, the seen/unseen calibration factor of Flora is kept at $0.75$. For PRIZM, we set $\lambda = 1$ in PPC (Eq.~\eqref{eq:signed}), $\beta = 0.3$, and $\gamma = 3$ in COA-Fuse (Eq.~\eqref{eq:adaptive}) before running any experiment. None of these values was tuned on unseen classes; Section~\ref{sec:beta} reports the sensitivity to $\beta$. All experiments use a fixed random seed (seed~1) unless stated otherwise (Section~\ref{sec:seeds} additionally reports seeds 2 and 3 for ZSL). Standard ZSL and ablation runs were executed on an NVIDIA L4 GPU, and GZSL, low-shot, and random-split runs on an NVIDIA RTX 5090 GPU; because Flora is non-deterministic across GPU types, every controlled rerun (\emph{No KG} and all ablation rows compared together) uses the same hardware and seed.

\paragraph{Evaluation protocol.}
Following Flora, the model is evaluated every 10 iterations during the flow-matching phase, which gives 20 checkpoints. We report the best unseen accuracy over these checkpoints in ZSL, and $\mathcal{S}$ and $\mathcal{U}$ at the checkpoint with the highest $\mathcal{H}$ in GZSL. Because this protocol selects among checkpoints evaluated on the test set, we also report the mean over the last ten checkpoints (last-10) where a comparison depends on it. Flora's own numbers are taken from its paper.

\subsection{Comparison with State-of-the-Art}

Table~\ref{tab:ntu-comparison} compares PRIZM with state-of-the-art ZSAR methods on NTU-60 and NTU-120. Flora numbers are the published ones; controlled comparisons with our own rerun of Flora without KG (\emph{No KG}) are given in Section~\ref{sec:ablation}.

\paragraph{Zero-shot learning.}
With STAR features, PRIZM obtains a four-split mean of 74.1\%, compared with 73.0\% reported for Flora and 73.1\% for our No KG rerun. Relative to the rerun, which isolates the effect of the KG, PRIZM improves 48/12 by 1.6 and 110/10 by 3.0 points, and leaves 55/5 ($-0.4$) and 96/24 ($-0.1$) essentially unchanged. The 96/24 result (67.0\%) is the best among the compared methods, but its margin over the published Flora number reflects run-to-run variation rather than the KG, since the No KG rerun already reaches 67.1\% on this split. On 55/5, almost all remaining errors are confusions between \emph{reading} and \emph{writing}, whose skeleton features overlap strongly, which leaves little room for a semantic prior. With SynSE features, PRIZM$^\dagger$ reaches 81.4\% on 110/10, the best result on this split and 1.8 points above Flora$^\dagger$. Its four-split mean is 74.5\%, compared with 74.4\% for Flora$^\dagger$, and on the other three splits it is within one point of Flora$^\dagger$. With both feature sets, the largest and most reliable gain is on 110/10.

\paragraph{Generalized zero-shot learning.}
With STAR features, PRIZM improves $\mathcal{H}$ over the published Flora results on 48/12 (+1.8), 110/10 (+0.7), and 96/24 (+2.6, best overall at 55.5\%), and is lower on 55/5 ($-0.4$). Its four-split mean $\mathcal{H}$ is 63.9\%, compared with 62.7\% for Flora. On NTU-60, where we also reran Flora without KG in GZSL, PRIZM is 3.1 points above the rerun on 48/12 ($\mathcal{H}$ 58.4\% vs.\ 55.3\%) and 0.2 points above it on 55/5. The gains on NTU-120 come mainly from the unseen classes: PRIZM and PRIZM$^\dagger$ obtain the two highest $\mathcal{U}$ among all methods on both 110/10 (72.9\% and 72.4\%) and 96/24 (61.6\% and 53.8\%). This is partly paid for by the seen classes, for example on 96/24, where $\mathcal{S}$ is 50.5\% compared with 53.7\% for Flora. With SynSE features, PRIZM$^\dagger$ obtains the best $\mathcal{H}$ on 110/10 (67.5\%, +1.4 over Flora$^\dagger$) and is on par with Flora$^\dagger$ on 96/24 (53.4\% vs.\ 53.2\%), but is lower on 55/5 ($-1.3$) and 48/12 ($-1.2$). Its four-split mean $\mathcal{H}$ is 64.1\%, compared with 64.3\% for Flora$^\dagger$. As in ZSL, 110/10 is the split where the improvement holds for both feature sets.

\paragraph{Low-shot training.}
Following Flora, we also train with only 1\% or 10\% of the training samples of each seen class, using SynSE features, and evaluate ZSL on the full unseen test set (Table~\ref{tab:lowshot}). The subset is drawn at random for each class with the fixed run seed and contains at least one sample per class. PRIZM obtains the best accuracy on 48/12 and 110/10 at both ratios and on 96/24 at 10\%, ties with Flora on 96/24 at 1\% (58.0\%), and is second to Flora on 55/5 at both ratios. Its four-split mean is 68.9\% at 1\% and 74.4\% at 10\%, compared with 66.2\% and 71.4\% for Flora. The largest margins are on 48/12 (+8.5 and +6.5 points) and 110/10 (+3.8 and +5.3 points). With 10\% of the data, PRIZM is close to its full-data result (74.4\% vs.\ 74.5\%), and with 1\% it loses 5.6 points on average, whereas Flora loses 8.2 points between its full-data and 1\% results. This is consistent with the KG acting as a fixed prior that does not have to be learned from the seen samples. These numbers, however, compare single-run results with the published Flora numbers, and at 1\%, where each seen class keeps only 5 to 6 training samples, the choice of the subset adds variance. On 55/5 at 1\%, for example, our own No KG rerun reaches 82.6\%, compared with 81.6\% for PRIZM. The low-shot comparison should therefore be read as indicative.

\paragraph{Random splits.}
Following Flora, we also evaluate on the random-split benchmark of SA-DVAE~\cite{savae24}, which provides three random seen/unseen splits per dataset and skeleton features extracted with ST-GCN, and report the mean over the three splits (Table~\ref{tab:random}). We use the configurations released by Flora for each split. Besides NTU-60 and NTU-120, this benchmark includes PKU-MMD I~\cite{pku} with 46 seen and 5 unseen classes. Its 51 classes all correspond to actions of NTU RGB+D 120, so the rules of 49 classes are copied unchanged from the corresponding NTU classes. The two remaining classes, \emph{put something inside pocket} and \emph{take out something from pocket}, both correspond to the single NTU class \emph{reach into pocket}; for them we added a hand-in-front primitive and opposite sequences, following the pattern of the other reversed pairs and using only the class names and the primitive vocabulary. ORE then separates all four reversed pairs of PKU-MMD, as on NTU. In ZSL, PRIZM obtains the best result on NTU-120 (72.9\%, +1.7 over Flora), is matches Flora on NTU-60 (88.6\%, below TDSM at 88.9\%), and is lower on PKU-MMD (69.4\% vs.\ 71.6\%). The splits of a dataset differ strongly in difficulty: on NTU-60, the three splits give 85.1\%, 95.0\%, and 85.6\%, and on PKU-MMD 51.3\%, 75.1\%, and 81.8\%. In GZSL, PRIZM obtains the best $\mathcal{H}$ on NTU-60 (81.4\%, +1.2 over Flora) and on NTU-120 (64.8\%, +1.8 over Flora; 71.5\%, 56.3\%, and 66.7\% on the three splits), and is lower than Flora and FS-VAE on PKU-MMD (58.1\% vs.\ 59.5\% and 59.0\%), where the three splits give 47.7\%, 59.5\%, and 67.2\%. As in the other benchmarks, the clearest gain is on NTU-120 110/10.

\begin{table}[htbp]
\centering
\caption{Performance comparison on NTU-60 (Xsub) and NTU-120 (Xsub). The best and second-best results in each column are marked in \textcolor{red}{red} and \textcolor{blue}{blue}, respectively. $\dagger$ denotes methods using SynSE-based~\cite{synse21} 4s-Shift-GCN features, while others use STAR-based~\cite{star24} 1s-Shift-GCN ones. $\ddagger$ indicates two-stream fusion; others are single-stream.}
\label{tab:ntu-comparison}
\scriptsize
\setlength{\tabcolsep}{2.3pt}
\renewcommand{\arraystretch}{1.1}
\newcommand{\best}[1]{\textcolor{red}{#1}}
\newcommand{\second}[1]{\textcolor{blue}{#1}}
\resizebox{\textwidth}{!}{%
\begin{tabular}{ll|cccc|cccc|cccc|cccc}
\hline
\multirow{3}{*}{Method} & \multirow{3}{*}{Venue}
& \multicolumn{8}{c|}{NTU RGB+D 60 (Xsub)}
& \multicolumn{8}{c}{NTU RGB+D 120 (Xsub)} \\
\cline{3-18}
& & \multicolumn{4}{c|}{55/5 Split}
& \multicolumn{4}{c|}{48/12 Split}
& \multicolumn{4}{c|}{110/10 Split}
& \multicolumn{4}{c}{96/24 Split} \\
\cline{3-18}
& & ZSL & \multicolumn{3}{c|}{GZSL}
& ZSL & \multicolumn{3}{c|}{GZSL}
& ZSL & \multicolumn{3}{c|}{GZSL}
& ZSL & \multicolumn{3}{c}{GZSL} \\
& & Acc & S & U & H & Acc & S & U & H & Acc & S & U & H & Acc & S & U & H \\
\hline
ReViSE~\cite{reVISE17} & ICCV 2017 & 69.5 & 40.8 & 50.2 & 45.0 & 24.0 & 21.8 & 14.8 & 17.6 & 19.8 & 0.6 & 14.5 & 1.1 & 8.5 & 3.4 & 1.5 & 2.1 \\
JPoSE~\cite{jpose19} & ICCV 2019 & 73.7 & 66.5 & 53.5 & 59.3 & 27.5 & 28.6 & 18.7 & 22.6 & 57.3 & 53.6 & 11.6 & 19.1 & 38.1 & 41.0 & 3.8 & 6.9 \\
CADA-VAE~\cite{cadavae19} & CVPR 2019 & 76.9 & 56.1 & 56.0 & 56.0 & 32.1 & 50.4 & 25.0 & 33.4 & 52.5 & 50.2 & 43.9 & 46.8 & 38.7 & 48.3 & 27.5 & 35.1 \\
SynSE~\cite{synse21} & ICIP 2021 & 71.9 & 51.3 & 47.4 & 49.2 & 31.3 & 44.1 & 22.9 & 30.1 & 52.4 & 57.3 & 43.2 & 49.5 & 41.9 & 48.1 & 32.9 & 39.1 \\
GZSSAR$^{\dagger}$~\cite{gzssar23} & IGIG 2023 & 83.6 & 71.7 & 66.2 & 68.8 & 49.2 & 58.8 & 40.0 & 47.6 & 71.2 & 46.8 & 68.3 & 55.6 & 59.7 & 56.8 & 48.6 & 52.4 \\
SMIE~\cite{smie23} & ACM MM 2023 & 77.9 & -- & -- & -- & 41.5 & -- & -- & -- & 61.3 & -- & -- & -- & 42.3 & -- & -- & -- \\
PURLS~\cite{purls24} & CVPR 2024 & 79.2 & -- & -- & -- & 41.0 & -- & -- & -- & 72.0 & -- & -- & -- & 52.0 & -- & -- & -- \\
SA-DVAE~\cite{savae24} & ECCV 2024 & 82.4 & 62.8 & 70.8 & 66.3 & 41.4 & 50.2 & 36.9 & 42.6 & 68.8 & 61.1 & 59.8 & 60.4 & 46.1 & \second{58.8} & 35.8 & 44.5 \\
STAR~\cite{star24} & ACM MM 2024 & 81.4 & 69.0 & 69.9 & 69.4 & 45.1 & 62.7 & 37.0 & 46.6 & 63.3 & 59.9 & 52.7 & 56.1 & 44.3 & 51.2 & 36.9 & 42.9 \\
STAR++~\cite{starpp26} & TCSVT 2026 & 84.4 & 61.1 & 73.6 & 66.8 & 49.5 & 58.2 & 40.4 & 47.7 & 72.0 & 59.0 & 55.4 & 57.2 & 53.5 & 52.8 & 45.2 & 48.7 \\
DVTA$^{\dagger}$~\cite{dvta25} & PR 2025 & 79.3 & -- & -- & -- & 44.1 & -- & -- & -- & 74.9 & -- & -- & -- & 51.8 & -- & -- & -- \\
InfoCPL$^{\dagger}$~\cite{infocpl25} & TMM 2025 & 85.9 & -- & -- & -- & 53.3 & -- & -- & -- & 74.8 & -- & -- & -- & 60.1 & -- & -- & -- \\
ScoPLe$^{\dagger}$~\cite{scople25} & CVPR 2025 & 84.1 & 69.6 & 71.9 & 70.8 & 53.0 & 54.5 & \best{61.8} & 57.9 & 74.5 & 63.5 & 61.1 & 62.3 & 52.2 & 53.3 & 51.2 & 52.2 \\
Neuron$^{\ddagger}$~\cite{neuron25} & CVPR 2025 & \best{86.9} & 69.1 & 73.8 & 71.4 & 62.7 & 61.6 & 56.8 & 59.1 & 71.5 & \best{67.6} & 59.5 & 63.3 & 57.1 & \best{67.5} & 44.4 & 53.6 \\
FS-VAE$^{\dagger}$~\cite{fsvae25} & ICCV 2025 & \best{86.9} & 77.0 & 74.5 & 75.7 & 57.2 & 56.2 & 48.6 & 52.1 & 74.4 & 59.2 & 67.9 & 63.3 & 62.5 & 57.8 & 51.9 & \second{54.7} \\
TDSM$^{\dagger}$~\cite{tdsm25} & ICCV 2025 & \second{86.5} & -- & -- & -- & 56.0 & -- & -- & -- & 74.2 & -- & -- & -- & 65.1 & -- & -- & -- \\
\textbf{Flora}~\cite{flora26} & CVPR 2026 & 85.8 & 77.7 & \second{75.6} & \second{76.6} & 61.5 & 66.9 & 49.0 & 56.6 & \second{80.7} & 59.8 & 70.5 & 64.7 & 64.1 & 53.7 & 52.2 & 52.9 \\
\textbf{Flora}$^{\dagger}$~\cite{flora26} & CVPR 2026 & 86.3 & 75.9 & \best{78.8} & \best{77.4} & \best{65.3} & 63.7 & \second{57.5} & \best{60.5} & 79.6 & \second{66.2} & 66.0 & \second{66.1} & 66.4 & 55.9 & 50.7 & 53.2 \\
\hline
\textbf{PRIZM (Ours)} & This work & 85.1 & \second{78.0} & 74.4 & 76.2 & 63.8 & \best{70.8} & 49.7 & 58.4 & 80.6 & 59.6 & \second{72.4} & 65.4 & \best{67.0} & 50.5 & \best{61.6} & \best{55.5} \\
\textbf{PRIZM (Ours)}$^{\dagger}$ & This work & 85.3 & \best{78.3} & 74.0 & 76.1 & \second{64.7} & \second{68.4} & 52.3 & \second{59.3} & \best{81.4} & 62.9 & \best{72.9} & \best{67.5} & \second{66.8} & 53.1 & \second{53.8} & 53.4 \\
\hline
\end{tabular}%
}
\end{table}

\begin{table}[htbp]
\centering
\caption{ZSL comparison under low-shot training (accuracy \%, SynSE-based~\cite{synse21} 4s-Shift-GCN features). 1\% and 10\% denote the fraction of training samples kept for each seen class. Results of the other methods are taken from Flora~\cite{flora26}. The best and second-best results in each column are marked in \textcolor{red}{red} and \textcolor{blue}{blue}, respectively.}
\label{tab:lowshot}
\small
\setlength{\tabcolsep}{4pt}
\newcommand{\best}[1]{\textcolor{red}{#1}}
\newcommand{\second}[1]{\textcolor{blue}{#1}}
\begin{tabular}{l|cc|cc|cc|cc}
\hline
\multirow{3}{*}{Method} & \multicolumn{4}{c|}{NTU RGB+D 60 (Xsub)} & \multicolumn{4}{c}{NTU RGB+D 120 (Xsub)} \\
\cline{2-9}
& \multicolumn{2}{c|}{55/5} & \multicolumn{2}{c|}{48/12} & \multicolumn{2}{c|}{110/10} & \multicolumn{2}{c}{96/24} \\
& 1\% & 10\% & 1\% & 10\% & 1\% & 10\% & 1\% & 10\% \\
\hline
ReViSE~\cite{reVISE17} & 51.0 & 58.0 & 9.8 & 15.6 & 14.8 & 23.6 & 5.2 & 7.8 \\
JPoSE~\cite{jpose19} & 33.0 & 62.1 & 23.8 & 28.3 & 15.6 & 49.9 & 8.6 & 33.9 \\
CADA-VAE~\cite{cadavae19} & 76.6 & 76.9 & 24.3 & 27.6 & 29.9 & 39.1 & 25.4 & 25.0 \\
SynSE~\cite{synse21} & 44.3 & 42.8 & 18.6 & 17.3 & 56.0 & 56.0 & 24.1 & 26.1 \\
SMIE~\cite{smie23} & 43.8 & 76.9 & 29.3 & 38.1 & 36.0 & 58.1 & 13.9 & 34.4 \\
SA-DVAE~\cite{savae24} & 21.1 & 60.4 & 18.7 & 20.0 & 14.4 & 40.9 & 9.5 & 21.9 \\
STAR~\cite{star24} & 40.6 & 77.0 & 11.6 & 35.3 & 18.9 & 46.8 & 8.6 & 33.0 \\
Neuron~\cite{neuron25} & 47.7 & 79.4 & 20.7 & 45.3 & 28.8 & 62.8 & 10.2 & 33.5 \\
FS-VAE~\cite{fsvae25} & 79.3 & 79.4 & 38.0 & 38.7 & 72.7 & 69.6 & \second{50.2} & 47.9 \\
TDSM~\cite{tdsm25} & 78.5 & 82.3 & 32.1 & 52.4 & 63.3 & 66.3 & 43.9 & 55.1 \\
\textbf{Flora}~\cite{flora26} & \best{82.8} & \best{85.6} & \second{46.5} & \second{56.1} & \second{77.4} & \second{78.1} & \best{58.0} & \second{65.9} \\
\hline
\textbf{PRIZM (Ours)} & \second{81.6} & \second{84.8} & \best{55.0} & \best{62.6} & \best{81.2} & \best{83.4} & \best{58.0} & \best{66.8} \\
\hline
\end{tabular}
\end{table}

\begin{table}[htbp]
\centering
\caption{Average performance over the three random seen/unseen splits of SA-DVAE~\cite{savae24} with ST-GCN features (ZSL accuracy and GZSL harmonic mean $\mathcal{H}$, \%). Results of the other methods are taken from Flora~\cite{flora26}. The best and second-best results in each column are marked in \textcolor{red}{red} and \textcolor{blue}{blue}, respectively.}
\label{tab:random}
\small
\setlength{\tabcolsep}{5pt}
\newcommand{\best}[1]{\textcolor{red}{#1}}
\newcommand{\second}[1]{\textcolor{blue}{#1}}
\begin{tabular}{l|cc|cc|cc}
\hline
\multirow{2}{*}{Method} & \multicolumn{2}{c|}{NTU-60 55/5} & \multicolumn{2}{c|}{NTU-120 110/10} & \multicolumn{2}{c}{PKU-MMD I 46/5} \\
& ZSL & GZSL & ZSL & GZSL & ZSL & GZSL \\
\hline
ReViSE~\cite{reVISE17} & 60.9 & 60.3 & 44.9 & 40.3 & 59.3 & 49.8 \\
JPoSE~\cite{jpose19} & 59.4 & 60.1 & 46.7 & 43.7 & 57.2 & 51.6 \\
CADA-VAE~\cite{cadavae19} & 61.8 & 66.4 & 45.2 & 45.6 & 60.7 & 45.8 \\
SynSE~\cite{synse21} & 64.2 & 67.5 & 47.3 & 43.5 & 60.8 & 49.5 \\
SMIE~\cite{smie23} & 65.1 & -- & 46.4 & -- & 60.8 & -- \\
SA-DVAE~\cite{savae24} & 84.2 & 75.3 & 50.7 & 47.5 & 66.5 & 54.7 \\
ScoPLe~\cite{scople25} & 83.7 & 77.7 & 53.3 & 54.1 & \second{71.4} & 54.9 \\
TDSM~\cite{tdsm25} & \best{88.9} & -- & 69.5 & -- & 70.8 & -- \\
FS-VAE~\cite{fsvae25} & -- & -- & -- & -- & 71.2 & \second{59.0} \\
\textbf{Flora}~\cite{flora26} & \second{88.6} & \second{80.2} & \second{71.2} & \second{63.0} & \best{71.6} & \best{59.5} \\
\hline
\textbf{PRIZM (Ours)} & \second{88.6} & \best{81.4} & \best{72.9} & \best{64.8} & 69.4 & 58.1 \\
\hline
\end{tabular}
\end{table}

\subsection{Ablation Study}
\label{sec:ablation}

\paragraph{Components in ZSL.}
Table~\ref{tab:components} adds the components of PRIZM one at a time with STAR features. Our No KG rerun matches the four-split mean of the Flora paper (73.1\% vs.\ 73.0\%), although individual splits differ by about 3 points (77.6\% vs.\ 80.7\% on 110/10 and 67.1\% vs.\ 64.1\% on 96/24), which shows the size of single-run variation. Adding the 82-primitive knowledge graph with uniform composition and raw addition (\emph{KG}) raises the mean by 0.9 points, with gains on 48/12, 110/10, and 96/24. Adding ORE and PPC leaves the mean almost unchanged (73.9\%): it improves 48/12 by 1.1 points but lowers 55/5 and 96/24. Its purpose is mainly representational, namely to separate reversed-order classes (Section~\ref{sec:further}), and in the standard splits this separation is not reflected in accuracy. Replacing raw addition by COA-Fuse gives the highest mean (74.1\%, +1.0 over No KG). This gain comes mainly from 110/10 (+1.4 over raw addition), while the other three splits change by less than 0.6 points.

\begin{table}[htbp]
\centering
\caption{Component ablation of PRIZM (ZSL accuracy \%, STAR features). Each row adds one component to the row above. \emph{KG} uses the 82 primitives with uniform composition and raw addition (without ORE, PPC, or COA-Fuse). \emph{No KG} is our controlled rerun of Flora.}
\label{tab:components}
\small
\begin{tabular}{l|cccc}
\hline
Semantics & 55/5 & 48/12 & 110/10 & 96/24 \\
\hline
No KG & \textbf{85.5} & 62.1 & 77.6 & 67.1 \\
\hline
KG & 85.4 & 62.9 & 79.2 & \textbf{68.5} \\
\quad + ORE + PPC & 84.7 & \textbf{64.0} & 79.2 & 67.5 \\
\quad + COA-Fuse (= PRIZM) & 85.1 & 63.8 & \textbf{80.6} & 67.0 \\
\hline
\end{tabular}
\end{table}

\paragraph{Steps of COA-Fuse.}
Table~\ref{tab:coa-steps} breaks down COA-Fuse into its three parameter-free operations on 48/12 and 110/10 (STAR features). Starting from ORE + PPC, centering (\emph{C-Fuse}) provides the main jump on 110/10 (+1.0 points, from 79.2\% to 80.2\%) by removing the common anisotropic component across classes. Adding orthogonalization (\emph{CO-Fuse}) achieves the highest accuracy on 48/12 (64.1\%) and maintains 80.2\% on 110/10. Adding the adaptive class weight (\emph{COA-Fuse}) achieves the highest accuracy on 110/10 (80.6\%, +1.4 over ORE + PPC), while 48/12 stays within 0.3 points of CO-Fuse.

\begin{table}[htbp]
\centering
\caption{Step-by-step ablation of COA-Fuse on NTU-60 48/12 and NTU-120 110/10 (ZSL accuracy \%, STAR features). \emph{ORE + PPC} and \emph{COA-Fuse} match Table~\ref{tab:components}.}
\label{tab:coa-steps}
\small
\begin{tabular}{l|cc}
\hline
Fusion & NTU-60 48/12 & NTU-120 110/10 \\
\hline
ORE + PPC & 64.0 & 79.2 \\
\quad + Center (C-Fuse) & 63.7 & 80.2 \\
\quad + Orth (CO-Fuse) & \textbf{64.1} & 80.2 \\
\quad + Adaptive (COA-Fuse = PRIZM) & 63.8 & \textbf{80.6} \\
\hline
\end{tabular}
\end{table}

\paragraph{Sensitivity to $\beta$.}
\label{sec:beta}
Table~\ref{tab:beta} varies the KG weight $\beta$ for C-Fuse, that is, COA-Fuse with the centering step only. Since ZSL provides no validation set of unseen classes, $\beta = 0.3$ was fixed before this sweep. On 48/12, accuracy stays at 63.7\% for $\beta \in \{0.1, 0.3\}$ and drops as $\beta$ increases further, reaching 58.7\% at $\beta = 1.0$ (3.4 points below No KG). On 110/10, accuracy peaks at $\beta = 0.5$ (81.6\%) and every tested value is above No KG. Across both splits, the single-seed mean peaks at the pre-selected $\beta = 0.3$ (71.9\%, compared with 71.3\%, 71.8\%, and 69.3\% for $\beta = 0.1$, $0.5$, and $1.0$). Averaged over three seeds, the mean is likewise highest at $\beta = 0.3$ (71.4\%, compared with 70.7\%, 71.1\%, and 69.3\%). The value fixed in advance therefore lies in the favorable range $[0.3, 0.5]$. The method is nevertheless sensitive to $\beta$ at the level of about one point, and a large KG weight can hurt.

\begin{table}[htbp]
\centering
\caption{Sensitivity to the KG weight $\beta$ (C-Fuse, ZSL accuracy \%, STAR features). $\beta = 0.3$ was fixed in advance for all main experiments; ZSL has no validation split on unseen classes.}
\label{tab:beta}
\small
\begin{tabular}{l|cccc|c}
\hline
$\beta$ & 0.1 & 0.3 & 0.5 & 1.0 & No KG \\
\hline
NTU-60 48/12 & \textbf{63.7} & \textbf{63.7} & 62.0 & 58.7 & 62.1 \\
NTU-120 110/10 & 78.9 & 80.2 & \textbf{81.6} & 79.8 & 77.6 \\
\hline
Mean & 71.3 & \textbf{71.9} & 71.8 & 69.3 & 69.9 \\
\hline
\end{tabular}
\end{table}

\paragraph{Variation across seeds.}
\label{sec:seeds}
Table~\ref{tab:ntu-comparison} reports seed~1. In ZSL, we also ran seeds 2 and 3. The four-split mean of PRIZM is 74.1\%, 74.5\%, and 74.1\% for seeds 1 to 3 with STAR features, and 74.5\%, 73.5\%, and 74.5\% with SynSE features, so the ranking of seeds is not consistent across feature sets. With STAR features, the three-seed means are $85.2 \pm 0.3$, $62.4 \pm 1.3$, $81.1 \pm 0.5$, and $68.2 \pm 1.2$ on 55/5, 48/12, 110/10, and 96/24 (mean 74.2\%), compared with 72.9\% for No KG and 74.0\% for the base KG. On 48/12 and 96/24, the standard deviation is as large as most differences between KG variants. Two effects exceed this variation: the improvement of the KG over No KG in the four-split mean (+1.0 for the base KG and +1.3 for PRIZM), and the improvement of COA-Fuse over the base KG on 110/10 (+1.4 on average, positive for all three seeds). The average difference between COA-Fuse and the base KG over the four splits (+0.3) is within seed variation. GZSL results are available for seed~1 only.

\subsection{Further Analyses}
\label{sec:further}

\paragraph{Shuffled-KG control.}
An additional semantic term could improve accuracy for reasons unrelated to its content, for example through the extra centering or the change of scale. To test this, we permute the KG rows across classes without fixed points, so that each class receives the rules of another class, and keep the primitive counts, the weight distributions, and all fusion operations unchanged (Table~\ref{tab:ablation-shuffle}). The permutation also changes the KG similarity used by the adaptive weight. On 48/12, the shuffled KG lowers accuracy by 6.1 to 7.0 points relative to the real KG and by 4.5 to 5.2 points below No KG across raw addition, C-Fuse, and COA-Fuse. On 110/10, the shuffled KG removes the whole gain: all shuffled variants (77.4\% to 77.8\%) lie within 0.2 points of No KG (77.6\%), whereas the real KG gives 79.2\% to 80.6\%. Over three seeds, the shuffled KG falls below No KG on both splits (by 2.4 to 2.9 points on 48/12 and by 0.5 to 0.7 points on 110/10), and the gap between the real and the shuffled KG (2.4 to 3.9 points) is about three times the standard deviation across seeds. The gains therefore depend on assigning the correct rules to each class, and the fusion operations do not produce them on their own.

\begin{table}[htbp]
\centering
\caption{Shuffled-KG control on NTU-60 48/12 and NTU-120 110/10 (ZSL accuracy \%, STAR features). \emph{Shuffled KG} permutes KG rows across classes with no fixed points while keeping all other components unchanged. A wrong KG falls well below \emph{No KG} on 48/12 and loses the whole gain of the real KG on 110/10.}
\label{tab:ablation-shuffle}
\small
\begin{tabular}{l|ccc|ccc}
\hline
\multirow{2}{*}{Semantics} & \multicolumn{3}{c|}{NTU-60 48/12} & \multicolumn{3}{c}{NTU-120 110/10} \\
& Real KG & Shuffled KG & $\Delta$ & Real KG & Shuffled KG & $\Delta$ \\
\hline
No KG & \multicolumn{3}{c|}{62.1} & \multicolumn{3}{c}{77.6} \\
\hline
ORE + PPC & \textbf{64.0} & 56.9 & +7.0 & 79.2 & 77.5 & +1.7 \\
C-Fuse & 63.7 & 57.6 & +6.1 & 80.2 & 77.4 & +2.8 \\
COA-Fuse (PRIZM) & 63.8 & 56.9 & +6.9 & \textbf{80.6} & 77.8 & +2.8 \\
\hline
\end{tabular}
\end{table}

\paragraph{KG-only semantics and primitive grounding.}
Training Flora on 48/12 with only the centered KG embeddings, without class names or LLM descriptions, gives 41.1\% accuracy, compared with 8.3\% for chance and 62.1\% for the text semantics of the No KG rerun. The KG alone thus carries substantial class information, although less than the text, and the two sources are complementary. To test whether the frozen skeleton features also contain primitive-level information, we train multi-label probes on the seen classes to predict the required primitives and evaluate them on unseen test samples. An MLP probe reaches an unseen AUC of 0.858 to 0.882 on the four splits, compared with 0.535 to 0.820 for probes trained on shuffled rules; the gap is large on 48/12, 110/10, and 96/24 and small on 55/5, where only 11 primitives can be evaluated. Primitive information is therefore present in the skeleton features. Classifying unseen actions from the predicted primitives alone, however, gives only 25.8\% to 72.9\% accuracy across splits, because many actions share most of their primitives. This explains why the KG is effective as a semantic prior but not as a standalone classifier.

\paragraph{Reversed-order class pairs.}
ORE removes all pairs of classes with identical rules (from seven to zero) and lowers the centered cosine similarity between the semantics of reversed-order pairs, such as \emph{wear jacket} and \emph{take off jacket}, from 0.99 to 0.78--0.81. In all four standard splits, however, each reversed-order pair has only one unseen member, while the other member is a seen class. Since ZSL evaluation is restricted to unseen labels, confusion within such a pair cannot occur, and we therefore claim only the representational separation reported above. Measuring its effect on accuracy would require a split in which both members of a pair are unseen.

\paragraph{Fixed fusion versus learned and decision-time alternatives.}
PRIZM keeps all KG operations parameter-free to avoid reasoning shortcuts~\cite{marconato26,andolfi25}. Consistent with this design, replacing the fixed composition with a learned module over the KG triples, an ESA module or an R-GCN, lowers accuracy on 48/12 below the fixed base KG, by 0.8 to 7.8 points for ESA depending on the regularization, because the learned parameters fit the seen classes. Using the KG on the skeleton side or at decision time does not help either. On 48/12, a learned logic head over predicted primitives, re-ranking with rules retrieved from the nearest seen samples, and re-ranking with the log-probabilities of the primitive probe all fall below PRIZM, by 4.2 to 17.4 points. The probe-based re-ranking improves the two NTU-120 splits (+0.8 and +1.1) but lowers both NTU-60 splits, so none of these variants gives a consistent gain. Finally, changing the phrasing of the primitive descriptions, either by joining the required primitives of a class into one sentence or by placing each primitive in a name-free context template, gives no consistent improvement over the plain descriptions. A likely reason is that a common template shifts all primitive embeddings in the same direction, which the centering step removes.

\bibliographystyle{ieeetr}
\bibliography{refs}

\end{document}
